\chapter{Background and Related Work}

The previous chapter introduced the project and its scope. This chapter provides the background needed to place that work in context. It begins with the learning problem that motivates the dissertation, then turns to interactive and immersive work in computer architecture and computer science education, and finally draws out the constraints and gaps that shaped the present project.

\section{Computer Architecture as a Learning Problem}

Computer architecture asks students to reason about a machine whose most important behaviour is hidden. A learner is expected to understand how instruction fields are decoded, how register values are selected, how the ALU is used for different classes of instruction, how memory is addressed, and how the program counter is updated. None of that is impossible to explain on paper, but it does place a heavy burden on the learner. Static diagrams show structure. Short traces show examples. The missing piece is often the transition between the two: how a flat picture becomes a moving process.

This problem is well known in practice. Architecture courses ask students to connect several kinds of representation at once: assembly syntax, instruction fields, datapath diagrams, register state, memory state, and control behaviour over time. Many students can solve a problem after seeing a worked example, but still struggle to explain why a branch changes the next program counter, why a load writes back a different value source from an arithmetic instruction, or why a sign-extended immediate matters at all. \citet{tlili2016improving} describe architecture as a subject that learners can find difficult and unmotivating under conventional classroom presentation. That diagnosis matches the teaching problem that motivated the current project.

The present dissertation focuses on a deliberately limited part of this wider area: instruction execution in a single-cycle MIPS datapath. That scope does not cover all of architecture, but it is large enough to expose the main difficulties that matter here. A learner still has to decode the instruction, identify source and destination registers, select operands, follow value flow through the datapath, and predict the resulting architectural state. The scope is therefore narrow enough to build and study within a Master's dissertation, while still being rich enough to raise the central educational questions.

\begin{figure}[htb]
\centering
\includegraphics[width=0.9\linewidth]{datapath single cycle.png}
\captionsetup{width=.9\linewidth}
\caption[Reference single-cycle MIPS datapath]{Datapath overview attributed to Jyotiprakash Mishra's blog post on the MIPS single-cycle datapath \citep{mishraMipsSingleCycleDatapath}.}
\label{fig:referenceSingleCycleDatapath}
\end{figure}

\section{Traditional Teaching and Its Limits}

The problem described in the previous section does not mean that conventional teaching is ineffective. In most programmes, computer architecture is taught progressively rather than all at once. Basic logic is introduced first, then gates, then combinations such as half adders and full adders, then larger datapath components, and only after that does the learner meet the processor as an integrated system. Trinity College Dublin's Computer Architecture I module reflects that pattern directly through its sequence of digital logic, register transfer language, ALU and shifter design, multiplexer and tristate buses, datapath design, and the instruction fetch-decode-execute cycle \citep{tcdComputerArchitectureI}. That kind of staged build-up matters, and the present dissertation does not argue against it.

Traditional tools also do several things very well. Lectures and tutorials give a structured order in which ideas can be introduced. Whiteboard derivations are good for slowing a problem down and showing how each step follows from the previous one. Static datapath diagrams are strong at showing structure, naming components, and fixing the shape of the machine in the learner's mind. Simulators and assembly tools are good for correctness, inspection, and repeated practice. They let students test programs, inspect registers and memory, and connect assembly syntax to machine behaviour in a more formal way than a verbal explanation alone.

These strengths are important because the present project depends on them. A learner cannot make sensible use of the VR environment without some prior exposure to the underlying architectural ideas. The VR prototype is not intended to replace the gradual conceptual build-up provided by lectures, diagrams, exercises, and conventional tools. It is intended to sit alongside them, especially at the point where a student has seen the structure of the datapath but still struggles to turn that structure into a working mental model of execution flow.

The limitation of conventional teaching is therefore narrower than a simple claim that the tools are ineffective. The issue is that some of the most useful traditional materials are strongest either at showing structure or at checking correctness. They are sometimes weaker at first-pass mental-model building for learners who need to see how information moves through the machine as a sequence of linked decisions. A datapath diagram can show where the ALU sits, and a simulator can show that a register changed, but for some students the connection between those two facts is not immediate. That gap is where a supplementary spatial tool may be useful.

\section{Interactive Support for Computer Architecture Learning}

Before focusing on VR, it is worth noting that architecture learning has already been explored through other interactive formats. \citet{tlili2016improving} present an educational mobile game for computer architecture and report that learners found the approach motivating and were willing to use it again. That work is not immersive, but it matters for two reasons. First, it suggests that the subject does respond well to interactive presentation. Second, it shows that the core issue is not VR by itself. The more basic issue is how to make abstract system behaviour easier to follow.

More directly related work appears in immersive and experiential systems. \citet{dascalu2017experiential} describe a VR environment for studying computer architecture through experiential learning. Their system asks the learner to assemble virtual computer components while receiving textual guidance. The emphasis is not identical to the present project, but the paper is still important because it treats computer architecture as a credible target for immersive learning rather than as a niche or novelty application. It also shows an early attempt to move architecture teaching away from passive explanation and toward structured activity inside a virtual environment.

\citet{tan2023vr} are closer still. Their VR serious game is explicitly aimed at the Computer Organisation and Architecture subject and focuses on CPU procedures that students often find hard to visualise in conventional teaching. They report positive learning efficiency together with an engaging experience. For the present dissertation, the main value of this paper is not merely that it uses VR, but that it treats architecture learning as a sequence of actions that can be staged and guided inside a game-like structure. That is close in spirit to the current prototype, even though the present work narrows the scope more tightly around instruction execution in a single-cycle datapath.

\citet{alnuaimi2025vr} study a VR environment for digital design and computer organisation laboratory learning. Their paper is useful partly because it remains close to the broader hardware-and-systems teaching space, and partly because it keeps usability in view throughout the evaluation. Their results point to a familiar lesson: even a promising educational VR system has to be learnable and comfortable enough to use before any learning claim can matter. For the present project, that is an important reminder. A tool intended to support understanding of hidden processes can just as easily create another layer of confusion if the interface is awkward or overloaded.

\section{Virtual Reality in Computer Science Education}

The broader computer science education literature gives a clearer sense of where architecture fits within immersive teaching work. \citet{pirker2020virtual} review VR in computer science education and note its promise for areas where abstract concepts benefit from interactive visualisation. \citet{agbo2021application} reach a similar conclusion from a larger bibliometric and content review, showing that VR has been explored across a range of computer science topics and that interest in the area has grown over time. \citet{janani2026virtual} provide the most recent scoping review and show that many VR applications in computer science education are directed toward conceptual learning rather than toward full practical programming tasks.

That pattern matters here. The present dissertation is not trying to teach learners how to write larger assembly programs or reason about the whole of digital logic design inside a headset. It is concerned with conceptual understanding of instruction execution. The wider reviews suggest that this is exactly the sort of problem for which VR is most often proposed: an abstract process, difficult to observe directly, that may become easier to follow when it is made spatial, interactive, and staged.

The more specific studies within that wider literature also point in the same direction. \citet{pirker2021potential} compare VR and desktop learning for sorting algorithms and report stronger subjective measures such as presence, absorption, and positive emotion in the VR condition. \citet{mohamad2023immersive} report positive results from a mobile VR application for learning Python coding. Neither paper is about architecture, but both are relevant because they treat VR as a way of making abstract computing concepts more concrete. They also show a recurring pattern in this literature: VR is often used less to replace existing formal tools than to provide a clearer first encounter with an otherwise difficult idea.

The same reviews also point to the limits of the field. \citet{janani2026virtual} note recurring usability issues, weak collaborative features, and a lack of clear instructional grounding in parts of the literature. \citet{agbo2021application} show that much of the work is still scattered across topics and methods. This is useful context for the present dissertation. It suggests that a contribution does not need to claim that VR in CS education is entirely new. The more realistic contribution is to show a careful application in a domain where the match between problem and medium is strong.

\section{Design Constraints and Transferable Lessons}

The case for VR is therefore not enough on its own. The next question is what kind of VR design is likely to help rather than hinder. \citet{cook2019challenges} discuss educational VR in terms of practical constraints as much as opportunity. Their work draws attention to issues such as deployment, accessibility, sustainability, and the amount of support institutions need before VR becomes a normal teaching tool. For a dissertation-scale prototype, this is a useful corrective. It shifts the emphasis away from spectacle and toward a more modest question: can a system be clear, usable, and worth the time required to use it?

\citet{korkut2023visualization} approach a related problem from the perspective of VR visualisation. They point out that immersive space does not automatically solve information-design problems. Poor layout, visual clutter, and badly organised interaction can still make a system hard to read. That lesson carries directly into architecture teaching, where the central challenge is already one of representing a hidden process without losing the learner in detail. If too much of the datapath is surfaced at once, VR may reproduce the original problem in another form.

\citet{samala2025virtual} make a similar point at the level of educational trends. Their review shows continuing interest in VR across education, but also repeated concerns about readiness, hardware access, discomfort, and the quality of actual learning design. This is helpful for keeping the present dissertation honest. A small architecture-learning prototype does not need to solve the institutional adoption problem, but it does need to avoid assuming that immersion alone will produce learning.

Some of the more transferable papers in the reading set sharpen this issue further by looking at abstract visualisation outside architecture. \citet{lee2020data} argue that VR can support deeper intuitive understanding when information is mapped in a way that users can feel and navigate, but they also warn that complexity can weaken focus. Work by Oberhauser \citep{oberhauser2018database,oberhauser2019database} and \citet{kellmann2024visualization} on database and linked-data visualisation is not about architecture, but it still matters because it shows that VR can be used to represent non-physical structures in ways that support navigation and comprehension. That lesson transfers well to the present project. A datapath is not a physical space that a learner can enter, but it can be turned into a navigable structure if the transformation is done carefully.

\section{Implications for the Present Project}

Several design implications follow from this literature. First, the case for using VR in computer architecture does not rest on novelty. It rests on a specific match between medium and problem: hidden, dynamic behaviour that learners often struggle to reconstruct from a static diagram. Second, the most useful systems are not simply more immersive versions of existing teaching material. They are selective. They decide what to externalise, what to keep simple, and where the learner should act rather than only observe.

Third, architecture learning seems to benefit from guidance and progression. This appears in game-based work such as Tlili et al. \citep{tlili2016improving} and Tan et al. \citep{tan2023vr}, and it also appears indirectly in the broader CS-education literature where prompts, panels, and staged activities are common design patterns \citep{janani2026virtual}. For the present dissertation, that supports the use of Learning, Practice, and Test as a progression rather than as three unrelated modes.

Fourth, the literature repeatedly warns against overload. A datapath can be made more visible in VR, but that does not mean every hidden detail should be turned into an explicit interaction. \citet{cook2019challenges}, \citet{korkut2023visualization}, \citet{samala2025virtual}, and \citet{lee2020data} all point in the same direction: if the learner is asked to manage too many objects, too much text, or too much movement at once, the value of the medium starts to disappear.

The gap addressed by the present project can therefore be stated fairly plainly. There is already evidence that interactive and immersive approaches can support computer architecture learning, and there is a broader body of work showing that VR can help with abstract computer science concepts. What remains less well served is a tightly scoped learning environment centred on manual instruction tracing through a single-cycle datapath, with explicit attention to operand selection, stage-by-stage decisions, and the final architectural state. That is the niche occupied by the current prototype. The next chapter explains how the design of that prototype was shaped by the literature reviewed here and by the teaching problem set out in Chapter 1.

\section{Summary}

The literature does not justify an argument that VR is automatically good for education, or even automatically good for computer architecture. It does justify a narrower claim. Computer architecture is a strong candidate for interactive support because much of the subject depends on understanding dynamic processes that are hard to see directly. Existing work shows that architecture learning has already been explored through mobile games, experiential environments, and VR serious games, while the wider VR-in-CS-education literature shows that conceptual learning is one of the medium's more plausible uses. The same literature also sets clear limits: guidance matters, usability matters, and information must be presented with restraint. Those are the conditions under which the present dissertation takes shape.
